\documentclass{article}
\usepackage{amsmath} % For math equations
\usepackage{amsfonts} % For math fonts
\usepackage{amssymb} % For math symbols
\usepackage{float}
\usepackage{enumitem}
\usepackage{graphicx}
\setlist[enumerate,1]{label=\arabic*.}
\setlist[enumerate,2]{label=\alph*.,itemindent=2em}

\title{Basic Exam Questions Practice}
\author{Asher Christian 006-150-286}
\date{}

\begin{document}

\section{F 15 Problem 7}
\emph{
    Let $A,B$ be two $4 \times 5$ matrices of rank $3$, and let $C = A^{T}B$. Find all possible values $r$ for the rank of $C$.
    If the rank $r$ is possible, find an explicit example of such matrices. Then prove that all other values are impossible
}
Available ranks
\begin{enumerate}
    \item rank 3
        \[
            B = \begin{pmatrix} 1 & 0 & 0 & 0 & 0 \\ 0 & 1 & 0 & 0 & 0 \\ 0 & 0 & 1 & 0 & 0 \\ 0 & 0 & 0 & 0 & 0 \end{pmatrix} \qquad A = B, A^{T} = \begin{pmatrix} 1 & 0 & 0 & 0 \\ 0 & 1 & 0 & 0 \\ 0 & 0 & 1 & 0 \\ 0 & 0 & 0 & 0 \\ 0 & 0 & 0 & 0 \end{pmatrix} 
        \] 
        \[
            A^{T}B = \begin{pmatrix} 1 & 0 & 0 & 0 & 0 \\ 0 & 1 & 0 & 0 & 0 \\ 0 & 0 & 1 & 0 & 0 \\ 0 & 0 & 0 & 0 & 0 \\ 0 & 0 & 0 & 0 & 0 \end{pmatrix} 
        \] 
        which has rank 3
    \item rank 2, keep $A$ as before but let
        \[
            B = \begin{pmatrix} 1 & 0 & 0 & 0 & 0 \\ 0 & 1 & 0 & 0 & 0 \\ 0 & 0 & 0 & 0 & 0 \\ 0 & 0 & 1 & 0 & 0 \end{pmatrix} 
        \] 
        then
        \[
            A^{T}B = \begin{pmatrix} 1 & 0 & 0 & 0 & 0 \\ 0 & 1 & 0 & 0 & 0 \\ 0 & 0 & 0 & 0 & 0 \\ 0 & 0 & 0 & 0 & 0 \\ 0 & 0 & 0 & 0 & 0  \end{pmatrix} 
        \] 
\end{enumerate}
I claim that these are the only two ranks possible of the resulting matrix. Indeed the rank of $A^{T}B$ must be less than or equal to 3
since the rank of $A^{T}$ is equal to the rank of $A$ by row-rank-column-rank equivalence and so $A^{T}X$ for any $X$ must have rank at most $3$. No
rank less than 2 is possible since $B$ can be thought of as a linear transformation from  $ \mathbb{R}^{5} \to \mathbb{R}^{4}$ with rank 3 which means It sends every vector into a 3 dimensional subsetof $ \mathbb{R}^{4}$ then $A^{T}$ is
a linear transformation from $ \mathbb{R}^{4} \to \mathbb{R}^{5}$ Of rank 3 so it sends all vectors in $ \mathbb{R}^{4} $ to a 3 dimensional subset of $ \mathbb{R}^{5}$. In the worst case 

\section{F 15 Problem 6}
\emph{
    Let $X := \mathbb{R} \setminus \{0\}$ Find a metric $ \rho $ on $X$ with the following properrties
    \begin{enumerate}
        \item $(X , \rho)$ is a complete metric space, and
        \item if $\{x_n\}_{n=1}^{\infty} \subset X$  and $x \in X$ then
             \[
                 \lim_{n\to \infty} |x_n - x| = 0 \iff x_n \rightarrow x \text{in} (X, \rho)
            \] 
    \end{enumerate}
    Prove both properties as well as all of your other assertions, in full detail
}
I first introduce the space
\[
    Y = ( \mathbb{R} \times \{0\}) \cup ( \mathbb{R} \times \{1\})
\] 
then I define a metric on $Y$, $D((x,i),(y,j))$ if  $i=j$ then  $D((x,i),(y,j)) = |x-y|$ the usual metric on  $ \mathbb{R}$
if $i =0, j=1$ then I define  $D((x,i),(y,j)) = |x| + |y-1|$ and if  $i=1, j=0$ then  $D((x,i),(y_j)) = |x-1| + |y|$. In essence
I am connecting these two lines i nthe metric so that  $(0,0)$ aligns with  $(1,1)$ ( they are the same point)
I propose the following metric
\[
\rho(x,y) =  \min(|x-y|, |x| + |y-1|, |x-1| + |y|)
\] 
now I need to prove that $\rho$ is a complete metric space.
We have the properties
\begin{enumerate}
    \item  $\rho(x,y) \ge 0$ for all $x,y$
    \item  $\rho(x,y) = 0 \implies x=y$ or  $x =0, y=1$ but  $0 \not\in X$ so this cannot be the case or $x =1, y=0$ which also cant happen so
         $\rho(x,y) = 0 \implies x=y$\
     \item 
         \[
         \rho(x,y) \le \rho(x,z) + \rho(z,y)
         \] 
\end{enumerate}

\section{F 15 Problem 5}
\emph{
    Let $F(x,y,z)$ be acontinuously differentiable function with non-vanishing partial
    derivatives at point $(0,0,0)$ define functions
    \[
    x = x(y,z), y = y(x,z), z = z(x,y)
    \] 
    as the solutions of the equation $F(x,y,z) = F(0,0,0)$ in the neighborhood of point  $(0,0)$ in the corresponding variables. Prove that
     \[
         \frac{\partial x}{\partial y} \frac{\partial y}{\partial z} \frac{\partial z}{\partial x} = -1
    \] 
    where the three partial derivates are taken at the point $(0,0)$ in the corresponding pair of variables.
}
consider
\[
    F'(0,0,0) = \begin{pmatrix} \frac{\partial F}{\partial x} & \frac{\partial F}{\partial y} & \frac{\partial F}{\partial z} \end{pmatrix} 
\] 
Now by the inverse function theorem 
\[
    x'(0,0) = (-A_x)^{-1}A_{y,z}
\] 
where
\[
    A = F'(0,0,0), A_x = \frac{\partial F}{\partial x}, A_{y,z} = \begin{pmatrix} \frac{\partial F}{\partial y} & \frac{\partial F}{\partial z} \end{pmatrix} 
\] 
so
\[
    \frac{\partial x}{\partial y} = ( - \frac{\partial F}{\partial x})^{-1}\frac{\partial F}{\partial y}
\] 
and
\[
    \frac{\partial x}{\partial y} \frac{\partial y}{\partial z} \frac{\partial z}{\partial x} = (- \frac{\partial F}{\partial x})^{-1} \frac{\partial F}{\partial y} (- \frac{\partial F}{\partial y})^{-1}\frac{ \partial F}{\partial z} (- \frac{\partial F}{\partial z})^{-1} \frac{ \partial F}{\partial x} = -1
\] 
the desired solution


\section{F 15 Problem 4}
\emph{
    Let $f_n : [0, \infty) \rightarrow \mathbb{R}$ be functions defined recursively by $f_1(x) :=0$
    and
    \[
        f_{n+1}(x) := e^{-2x} + \int_{0}^{x}f_n(t)e^{-2t}dt, \;\;\; n\ge 1
    \] 
    Show that $f(x) := \lim_{n\to \infty}f_n(x)$ exists for all $x \ge 0$ and identify  $f$ explicitely
}
Consider
\[
    g_n(x) = f_n(x) - f_{n-1}(x) = e^{-2x} + \int_{0}^{x}f_n(t)e^{-2t}dt - e^{-2x} - \int_{0}^{x}f_{n-1}(t)e^{-2t}dt = \int_{0}^{x}(f_n(t)-f_{n-1}(t))e^{-2t}dt
\] 
\[
    = \int_{0}^{x}g_{n-1}(t)e^{-2t}dt
\] 
so for fixed $x \ge 0$ using the supremum norm
\[
    ||g_n||_{[0,x]} \le ||g_{n-1}||_{[0,x]}\int_{0}^{x}e^{-2t}dt = ||g_{n-1}|| \frac{1-e^{-2x}}{2}
\] 
\[
    ||g_n||_{[0,x]} \le ||g_2||(\frac{1-e^{-2x}}{2})^{n-2}
\] 
and $||g_2||$ is finite
so
 \[
     ||g_n||_{[0,x]} \rightarrow 0
\] 
as $n \rightarrow \infty$ for all $x$ so
 \[
\sum_{k \ge 1}^{}g_k
\] 
is unfiormly convergent on $[0,x]$ so that  $f_n$ are cachy and the limit $f$ exists for all $x$.
The solution must satisfy
\[
f(x) = e^{-2x} + \int_{0}^{x}f(t)e^{-2t}dt
\] 
solving this ODE we get
\[
f(x) = 2 - e^{(\frac{1}{2}(1-e^{-2x}))}
\] 

\section{F 15 Problem 3}
\emph{
    Let $\{f_n\}$ be a sequence of continuous functions  $f_n :[-1,1] \rightarrow [0,1]$ such that, for each $x \in [-1,1]$
     \begin{enumerate}
         \item the sequence of numbers $\{f_n(x)\}_{n=1}^{\infty}$ is non-increasing, and
         \item $\lim_{n\to \infty}f_n(x)= 0$
    \end{enumerate}
    define
    \[
    g_n(x) := \sum_{m=1}^{n}(-1)^{m}f_m(x)
    \] 
    prove that $g_n(x)$ converges to some  $g(x) \in \mathbb{R}$ for each $x \in [-1,1]$
    and that the function  $g: [-1,1] \rightarrow \mathbb{R}$ is continuous on $[-1,1]$
}
First to show that the limit exists for each $x$ fix  $\epsilon > 0$ want to show there exists $N$ such that 
for all  $n, m> N, |g_n(x)- g_m(x)| < \epsilon$ for $m > n$
\[
g_m(x) - g_n(x) =  \sum_{i=n+1}^{m}(-1)^{i}f_i(x)
\] 
Let $N$ be such that $f_i(x) < \epsilon$ for all $ i > N$ then I assert that
\[
|\sum_{i=n+1}^{m}(-1)^{i}f_i(x)| < \epsilon
\] 
assume wlog that $n+1$ is even then we have that our sequence is bounded above by  $f_{n+1}(x)$ since the first term is bounded above by  $f_{n+1}(x)$ and the
sequence with  $f_{n+1}(x) - f_{n+2}(x) < f_{n+1}(x)$ by positivity and if we are adding a positive term then the positive term is less than the negative term from the previous part so the
term is less than the sum discarding the term and the previous term which inductively is less than  $f_{n+1}(x)$, more specifically if 
 \[
     \sum_{i=n+1}^{n+2m + 1}(-1)^{i}f_i(x) \le f_{n+1}(x) \implies \sum_{i=n+1}^{n+2m+2}(-1)^{i}f_i(x) \le \sum_{i=n+1}^{n+2m+3}(-1)^{i}f_i(x) \le f_{n+1}(x)
\] 
so the sequence $g_n(x)$ is cauchy for fixed  $x$ and so the limit is well defined. Now if $g_n$ converges to $g$ uniformly then we have that $g$ is continuous
this is because the uniform limit of continuous functions is continuous and since each $g_m$ is a finite sum of continuous functions it is continuous. Indeed consider
\[
    |g_n(x) - g(x)| = |\sum_{i=n+1}^{\infty}(-1)^{i}f_i(x)| \le f_{n+1}(x) \le \sup_{[-1,1]}f_{n+1}(x)
\] 
so it suffices to show that $\sup_{x}f_n(x) \rightarrow 0$ as $n \rightarrow \infty$. Assume not then there exists $\epsilon >0$ such that
for all $n$ there exists  $x_n$ such that  $f_n(x_n) \ge \epsilon$ then since $[-1,1]$ is compact 
there exists $(x_{n_k})_k$ such that  $(x_{n_k})_k \rightarrow y \in [-1,1]$ converges  by sequentially compact.
so 
\[
    f_{n_k}(x_{n_k}) \ge \epsilon \forall k
\] 
there exists $K$ such that  $f_k(y) < \frac{\epsilon}{3}$ for all $k \ge K$ and there exists  $\delta > 0$ so that $f_K(x) < \frac{2\epsilon}{3}$ forall $|x - y| < \delta$ and so by monotonicity
$f_k(x) < \frac{2\epsilon}{3}$ for all $k \ge K$ and $|x-y| < \delta$ but there exists some $k \ge K$ such that $|x_{n_k} - y| < \delta$ implies that $f_{n_k}(x_{n_k}) \le \frac{2\epsilon}{3}$ which is a contradiction
so it must be that $\sup_x(f_n(x)) \rightarrow 0$ as $n \rightarrow \infty$ and similarly that $g_n \rightarrow g$ uniformly which further implies that $g$ is continuous

\section{F 15 Problem 2}
\emph{
    Let $a,b \in \mathbb{R}$ obey $a < b$. Show that if  $g, h: [a,b] \rightarrow \mathbb{R}$ are 
    continuous with $h \ge 0$ then there is  $c \in [a,b]$ such that
     \[
    \int_{a}^{b}g(x)h(x)dx = g(c)\int_{a}^{b}h(x)dx
    \] 
}
    I want to use intermediate value theorem on the function
    \[
    g(x)\int_{a}^{b}h(t)dt
    \] 
    since $g$ is continuous on $[a,b]$ (compact) there exists  $x_1$ such that $g(x_1) \le g(x) \forall x \in [a,b]$ 
    and $x_2$ such that $g(x_2) \ge g(x) \forall x \in [a,b]$ and we have
    \[
    g(x_1)\int_{a}^{b}h(t)dt = \int_{a}^{b}g(x_1)h(t)dt \le \int_{a}^{b}g(t)h(t)dt
    \] 
    and
    \[
    g(x_2)\int_{a}^{b}h(t)dt = \int_{a}^{b}g(x_2)h(t)dt \ge \int_{a}^{b}g(t)h(t)dt
    \] 
    where the inequality comes from the fact that if $a < b$ then  $ah(t) \le bh(t)$ since  $h$ is strictly nonnegative.
    Thus $g(x)\int_{a}^{b}h(t)dt$ is less than and greater than $\int_{a}^{b}g(x)h(x)dx$ for some $x_1,x_2 \in [a,b]$ and so by intermediate value theorem
    of continuous functions there exists $c$ such that
    \[
    g(c)\int_{a}^{b}h(t)dt = \int_{a}^{b}h(x)g(x)dx
    \] 


\section{F 15 Problem 1}
\emph{
    Let $\{a_n\}_{n=1}^{\infty}$ be a sequence of positive numbers such that
    \[
        a_{n+m} \le a_n + a_m, \;\;\; m,n \ge 1
    \] 
    Prove that $\lim_{n\to \infty} \frac{a_n}{n}$ exists by showing
    \[
        \lim_{n\to \infty} \frac{a_n}{n} = \inf_{n \ge 1}\frac{a_n}{n}
    \] 
}
\\
We have that the limit exists iff $\limsup_{n\to \infty}\frac{a_n}{n} = \liminf_{n\to \infty}\frac{a_n}{n}$.
\[
    \limsup_{n\to \infty}\frac{a_n}{n} = \inf_{n \ge 1}\sup_{m \ge n}\frac{a_m}{m}
\] 
then consider
\[
    \sup_{m \ge n} \frac{a_m}{m}
\] 
we have
\[
m = qn + r
\] 
for some $q,r \in \mathbb{N}$, $r < n$. then we can say
\[
a_m \le q a_n + a_r
\] 
and
\[
    \frac{a_m}{m} \le q \frac{a_n}{m} + \frac{a_r}{m} \le \frac{a_n}{n} + \frac{a_r}{m}
\] 
since
\[
\frac{q}{m} \le \frac{1}{n}
\] 
and since $a_r$ is bounded since it can only range over the first $n$ elements that term must go to zero 
so we clearly have
\[
\limsup_{m\to \infty}\frac{a_m}{m} \le \frac{a_n}{n}
\] 
for all $n$ so 
\[
    \limsup_{n\to \infty}\frac{a_n}{n} \le \inf_{n \ge 1} \frac{a_n}{n}
\] 
additionally
\[
    \liminf_{n\to \infty}\frac{a_n}{n} = \sup_{n\ge 1}\inf_{m \ge n} \frac{a_m}{m} \ge \inf_{m \ge n}\frac{a_m}{m} \ge \inf_{m \ge 1} \frac{a_m}{m}
\] 
so we have
\[
    \limsup_{n\to \infty}\frac{a_n}{n} \le \inf_{n \ge 1} \frac{a_n}{n} \le \liminf_{n\to \infty} \frac{a_n}{n}
\] 
and thus the limsup is equal to the liminf which are together equal to the limit.

\end{document}
